\documentclass[12pt]{article}
\usepackage[utf8]{inputenc}
\usepackage{amsmath}    
\usepackage{amssymb}    
\usepackage{geometry}   
\geometry{a4paper, margin=1in}

\begin{document}

\section{Phase Accumulation and the Continuous Phase Model}

\subsection{Definition of Terms and Variables}
To analyze the conversion of digital bits into a Radio Frequency (RF) waveform, the following parameters are defined:
\begin{itemize}
    \item \textbf{Carrier Frequency ($f_c$):} The center frequency (e.g., 2.402 GHz).
    \item \textbf{Modulation Index ($h$):} A ratio defining the frequency swing relative to the bit rate. For BLE, $h = 0.5$.
    \item \textbf{Bit Duration ($T_b$):} The time allotted for one bit ($1 \, \mu s$ for 1 Mbps PHY).
    \item \textbf{Instantaneous Phase ($\theta(t)$):} The actual value used to drive the Sine/Cosine generator.
\end{itemize}

\subsection{Derivation of Frequency Deviation ($\Delta f$)}
The frequency deviation is the maximum distance the signal moves from the center carrier. In GFSK, this is derived from the modulation index $h$. 
\begin{equation}
h = \text{Total Frequency Swing} \times T_b = (2\Delta f) \cdot T_b
\end{equation}
Isolating $\Delta f$ through algebraic division:
\begin{equation}
\Delta f = \frac{h}{2T_b}
\end{equation}
Substituting the BLE standard values:
\begin{equation}
\Delta f = \frac{0.5}{2 \times 10^{-6}} = 250,000 \, \text{Hz} = 250 \, \text{kHz}
\end{equation}

\subsection{Step-by-Step Phase Integral Derivation}
The goal of the Baseband Processor is to generate a phase $\theta(t)$ that is the integral of the frequency. Given a sequence of bits $a_i$ and the shaping pulse $g(t)$, the instantaneous frequency deviation is:
\begin{equation}
f_{\text{dev}}(\tau) = \sum_{i} a_i g(\tau - iT_b)
\end{equation}
The phase is the integral of this frequency deviation over time:
\begin{equation}
\theta(t) = 2\pi h \int_{-\infty}^{t} \left[ \sum_{i} a_i g(\tau - iT_b) \right] d\tau
\end{equation}
Using the linearity property of integration, the summation is moved outside the integral:
\begin{equation}
\theta(t) = 2\pi h \sum_{i} a_i \left( \int_{-\infty}^{t} g(\tau - iT_b) d\tau \right)
\end{equation}
This result shows that the total phase is the cumulative sum of the areas under the shaped pulses.

\subsection{RTL Significance: From Calculus to Hardware}
In the RTL implementation of a BLE Baseband, this mathematical model is mapped to a \textbf{Phase Accumulator} within a Numerically Controlled Oscillator (NCO).
\begin{enumerate}
    \item \textbf{Continuous vs. Discrete:} The integral is replaced by a digital summation: $\theta[n] = \sum f_{word}[n]$.
    \item \textbf{Phase Continuity:} The accumulator ensures that when a bit changes from $+1$ to $-1$, the phase starts from its last value rather than resetting to zero. This prevents "Phase Jumps."
    \item \textbf{Hardware Efficiency:} The integral of the shaping pulse $\int g(\tau) d\tau$ is pre-calculated and stored in a Look-Up Table (LUT), reducing the real-time computational load on the FPGA/ASIC to simple addition.
\end{enumerate}

\end{document}